% Chapter 2 -- mirrors PDF structure and content (condensed where noted)
\chapter{Background Knowledge and Literature Review}

\section{Introduction}
The literature review focuses on existing research related to robotic exoskeletons for back support, particularly in the context of manual material handling (MMH) and load transportation tasks. Key themes include exoskeleton designs (passive vs.\ active, rigid vs.\ flexible), actuation systems (with emphasis on cable-driven and series elastic actuators), control strategies (including intention detection, sensor fusion, and machine learning), and personalization mechanisms. The review draws from systematic reviews, empirical studies, and technical papers to identify advancements, limitations, and gaps relevant to the project's goal of developing a software-centric enhancement module for a lower-limb exoskeleton. To provide a comprehensive overview, this review incorporates recent systematic analyses from 2020--2025, highlighting post-pandemic advancements in hybrid designs, AI integration, and industrial evaluations. For instance, Pesenti et al.~\cite{pesenti2021} and subsequent updates in Kermavnar et al.~\cite{kermavnar2021} and later works in bioengineering reviews classify over 30 back-support exoskeletons, emphasizing the shift toward soft and quasi-passive systems for improved user acceptance.

\section{Overview of Back-Support Exoskeletons for MMH}
Back-support exoskeletons (BSEs) are designed to mitigate low back pain and musculoskeletal disorders (MSDs) associated with repetitive lifting, bending, and load carrying in industries such as construction, logistics, and manufacturing. A systematic review by De Looze et al.~\cite{delooze2016} and updated analyses by Pesenti et al.~\cite{pesenti2021} classify BSEs into passive and active types. Passive BSEs, such as the Laevo or SPEXOR, rely on mechanical springs or elastic bands to store and release energy during trunk flexion-extension (TFE) motions, reducing erector spinae muscle activity by 10--40\% without requiring power sources. However, they offer limited adaptability to dynamic tasks and may impose resistance during non-assisted movements.

Recent evaluations, such as Luger et al.~\cite{luger2023}, evaluated the Laevo V2 passive exoskeleton in male workers and reported 12--19\% lower erector spinae activity during symmetric and asymmetric static holding tasks. Koopman et al.~\cite{koopman2020} further quantified that although passive exoskeletons effectively reduce peak L5/S1 moments during symmetric lifting, the moment reduction diminishes significantly when lifting involves trunk rotation or asymmetry---reinforcing the need for active, adaptive systems in real-world industrial settings where perfectly symmetric lifts are rare.

Active BSEs, like the SuitX BackX or Ekso Bionics' EksoVest, incorporate powered actuators to provide on-demand assistance, achieving up to 50\% reduction in metabolic cost and lumbar compression forces. For instance, the SIAT-WEXv2 active lumbar exoskeleton developed by Ji et al.~\cite{ji2020} reduced peak lumbar moments by 25--35\% during repetitive box lifting of 15--25~kg loads, confirming the effectiveness of combined hip and knee actuation in scenarios directly comparable to the $>20$~kg transportation tasks targeted in this project. Golabchi et al.~\cite{golabchi2023} highlight their potential in construction, where workers handle loads exceeding 20~kg, but note challenges in portability, weight (often $>5$~kg), and user comfort. A comprehensive review by Kermavnar et al.~\cite{kermavnar2021} of 28 BSEs for MMH tasks emphasizes that while active systems excel in heavy-load scenarios, they often suffer from rigid structures that restrict natural range of motion (RoM) in sagittal, frontal, and transverse planes, leading to misalignment and discomfort. Simon et al.~\cite{simon2021} observed slight forward tilt of the pelvis and reduced lumbar flexion when using a lift-assistive exoskeleton, underscoring the importance of compliant actuation designs that preserve physiological joint trajectories.

Long-term field evaluations show that BSE adoption increases with user acceptance, but dropout rates are high due to fit issues and lack of personalization. In load transportation contexts, studies on the HeroWear Apex exosuit indicate reductions in back forces during lifting and bending, with metabolic cost decreases of 7.9--23\% but highlight trade-offs in comfort for overhead tasks~\cite{yandell2020,dai2021}. Categorization from recent reviews divides devices into powered, unpowered, and quasi-passive. Powered types, such as the Cray X, reduce erector spinae activity by 16--37.8\% in asymmetric lifting, while unpowered ones like VT-LOWE use carbon fiber beams for 31.5\% peak muscle reduction. Quasi-passive hybrids combine springs with minimal actuators for 36.92\% activity reduction, offering energy efficiency.

The SPEXOR passive spinal exoskeleton reduced net metabolic cost by 18\% (range 11--25\%) during symmetric repetitive lifting~\cite{baltrusch2020}. More recent quasi-passive designs with compact variable gravity compensation modules have achieved comparable reductions in erector spinae activity with improved wearer comfort~\cite{song2024}. Long-term field adoption remains constrained by limited adaptability to highly dynamic or asymmetric tasks and by user-acceptance challenges~\cite{babic2021,luger2023}.

\section{Actuation Systems: Focus on Cable-Driven and Series Elastic Actuators}
Actuation design is critical for efficient torque delivery while ensuring safety and ergonomics. Traditional rigid actuators transmit force via linkages but can cause parasitic forces and limit RoM. Cable-driven systems address this by using flexible transmissions, allowing off-board or on-board placement for portability. Chen et al.~\cite{chen2019} and Li et al.~\cite{li2022} review and demonstrate cable-driven exoskeletons, noting their advantages in reducing structural weight and enabling multi-DoF compensation for misalignment. Recent advancements include Li et al.'s cable-driven asymmetric back exosuit, providing 38\% muscle activity reduction in lifting, and Ding et al.'s spring-cable-differential for passive assistance~\cite{ding2023}.

Series elastic actuators (SEAs) integrate compliant elements to enhance shock tolerance, force control, and human--exoskeleton interaction safety. Pratt and Williamson~\cite{pratt1995} pioneered SEAs, and recent applications in BSEs include Hyun et al.~\cite{hyun2020}. Cable-driven series elastic actuation (CSEA) proposed by Liao et al.~\cite{liao2024} uses a torsion spring--support beam mechanism with deflection constraints to operate in multiple modes: SEA mode for compliant, low-torque assistance and conventional stiff actuator (CSA) mode for high-torque output ($>50$~N$\cdot$m). Bench and human tests validated a 20--30\% reduction in erector spinae EMG during TFE, with stable torque control despite status transitions.

Poliero et al.~\cite{poliero2022} compared active and passive BSEs in dynamic tasks. Ding et al.~\cite{ding2024} introduced a differential SEA with cable routing for precise force control. In soft exoskeletons, twisted string actuators enable lightweight designs with high force density~\cite{xiloyannis2019}. Comparative analysis shows that cable-driven SEAs excel in flexibility but require careful routing and low-friction sheathing~\cite{poliero2022,liao2024}. Bowden cable routing and pre-tensioning, as in upper-limb systems~\cite{tsai2019}, minimize friction and hysteresis. For load transportation, these systems typically deliver 20--50~N$\cdot$m of torque---sufficient for assisting with loads greater than 20~kg---although challenges in battery life and heat dissipation remain.

\section{Control Strategies: Intention Detection, Sensor Fusion, and Machine Learning}
Effective BSE control requires real-time detection of user intent and environmental adaptation. Traditional methods use EMG or IMUs for gait event detection, but they struggle with transitions. Machine learning has emerged for intention prediction: Novak and Riener~\cite{novak2015} survey sensor fusion methods in wearable robotics and highlight intention detection as a central challenge. Separately, locomotion mode classification studies often use SVMs or LSTM networks and report $>90$\% accuracy in controlled settings. Recent applications include Jiang et al.'s EMG-based loading recognition for lumbar exoskeletons~\cite{jiang2024}. Chen et al.~\cite{chen2018} showed that simple kinematic thresholds on hip flexion angle and angular velocity can reliably detect the onset of lifting within 150--200~ms, offering a lightweight, non-ML baseline.

Sensor fusion, often via Kalman filters, integrates IMU data for accurate joint kinematics estimation with low latency ($<200$~ms). Lazzaroni et al.~\cite{lazzaroni2022} improved efficacy with accelerometer-based control during manual handling. Toxiri et al.~\cite{toxiri2018} use finite state machines to modulate assistance based on trunk angle, reducing user effort by 25\%. Personalization is a growing focus: generic assistance profiles underperform across diverse users~\cite{kermavnar2021,luger2023}. Reinforcement learning for control optimization can reduce metabolic cost by 13--24\% in walking/stairs~\cite{luo2023}. Xiang et al.~\cite{xiang2023} dynamically evaluated multiple back-support exoskeletons and reported 19--69\% reductions in peak back-muscle activity across manual handling tasks, highlighting needs for ML in asymmetric tasks.

\section{Evaluation Studies and User Acceptance}
Empirical studies provide quantitative insights. For passive BSEs, van Sluijs et al.~\cite{vanSluijs2023} reported physiological benefits in lifting/leaning, while Alemi et al.~\cite{alemi2020} reported 4--13\% reductions in energy expenditure during repetitive lifting. Qu et al.~\cite{qu2021} evaluated a commercial passive exoskeleton in real industrial lifting tasks and reported 18--32\% lower back-muscle activity (iEMG) and improved subjective comfort. Active systems reduce activity during lifting with subjective comfort improvements~\cite{schwartz2023}. Quasi-passive systems use Gaussian mixture models for periodic assistance~\cite{jamsek2020}. User acceptance studies show soft actuated BSEs influence multimodal measurements positively, but gender differences exist in overhead tasks~\cite{refai2023,nussbaum2023}.

\section{Gaps and Implications for the Project}
Tr\"oster et al.~\cite{troster2022} used a musculoskeletal model to compare stoop and squat lifting with and without back-support exoskeletons, concluding that assistance magnitude must be adjusted dynamically according to lifting technique---further motivating adaptive control. Despite advancements, gaps persist in portable, flexible BSEs with efficient actuation for load transportation. ML and sensor fusion show promise for proactive assistance, but few studies combine them with personalization for real-time adaptation~\cite{lin2025}. Lin et al.~\cite{lin2025} further highlight technical challenges in locomotion-assistance prototypes. The interim project plan positioned the work as enhancing an existing lower-limb prototype with MATLAB-based IMU fusion, SVM/LSTM-style intention detection, and trial-based fine-tuning, targeting $>90$\% classification accuracy and substantial (\(\sim\)30--50\%) reductions in estimated musculoskeletal strain under controlled validation---with the understanding that the LSTM path and strain metrics would follow incremental validation.

The present work addresses software-centric rapid prototyping on IMU data with MATLAB-based fusion and classification~\cite{zhang2012,chereshnev2018}, targeting $>90$\% accuracy on public gait data as a stepping stone toward hardware integration.

The present work directly addresses gaps repeatedly identified in recent reviews~\cite{pesenti2021,kermavnar2021,troster2022,lin2025}:
\begin{itemize}[noitemsep]
  \item Lack of real-time personalization and user-specific adaptation (future work).
  \item Insufficient handling of asymmetric and highly dynamic load-transportation tasks.
  \item Delayed or reactive assistance onset in active/quasi-passive systems.
  \item Limited transferability and rapid prototyping of advanced control on existing hardware.
\end{itemize}

\section{Summary Table of Representative Back-Support Exoskeletons}
\begin{table}[htbp]
  \centering
  \small
  \caption{Representative back-support exoskeletons (adapted from the literature).}
  \label{tab:exo_summary}
  \begin{tabular}{@{}p{2.6cm}p{1.2cm}p{1.1cm}p{3.2cm}p{3.0cm}@{}}
    \toprule
    \textbf{Device / Study} & \textbf{Type} & \textbf{Torque (N$\cdot$m)} & \textbf{Reported benefit} & \textbf{Key limitation} \\
    \midrule
    Laevo V2 (2023) & Passive & $\sim$35--40 & 12--19\% erector spinae EMG (static) & Resistance in extension; poor asymmetry \\
    SPEXOR (2020) & Passive & $\sim$40 & 11--25\% net metabolic cost & Fixed profile; added hip stiffness \\
    HeroWear Apex & Soft quasi-passive & 25--35 & 7.9--23\% metabolic; $\sim$20--30\% back force & Lower torque; limited overhead \\
    SuitX BackX & Active & 30--60 & Up to 50\% lumbar compression & Weight $>7$~kg; rigid RoM limits \\
    Cray X (2022) & Active & 50--80 & 16--37.8\% erector spinae EMG & High power; bulky battery \\
    SIAT-WEXv2 & Active hip+knee & $\sim$60 & 25--35\% peak lumbar moment (15--25~kg) & Complex mechanics; $>5$~kg \\
    VT-LOWE & Passive & $\sim$45 & 31.5\% peak back-muscle activity & No dynamic adaptation \\
    Liao et al.\ CSEA (2024) & Hybrid & $>50$ (CSA) & 20--30\% erector spinae EMG & Mode-switching complexity \\
    Auxivo LiftSuit (2024) & Quasi-passive & 30--45 & 15--28\% metabolic; $\sim$35\% EMG & Limited customization \\
    Hyundai / German Bionic & Active & 50--75 & 20--40\% back load reduction & Cost; rigid harness \\
    \bottomrule
  \end{tabular}
\end{table}
